\documentclass[UTF8]{ctexart} 
%\usepackage{homework}
\usepackage{geometry}
\geometry{a4paper,scale=0.8}
\usepackage[table,xcdraw]{xcolor}
\usepackage{graphicx} %插入图片的宏包
\usepackage{float} %设置图片浮动位置的宏包
\usepackage{subfigure} %插入多图时用子图显示的宏包
\usepackage{enumerate}
%\usepackage{fancyhdr}  % header，footer的设置
%\usepackage{extramarks}
\usepackage{amsmath}  % 数学公式
\usepackage{amsthm}
\usepackage{amsfonts}
\usepackage{tikz}  % 绘图
\usepackage{algorithm}  % 算法
\usepackage{algorithmicx}
\usepackage{algpseudocode}  % 伪代码
%\usepackage[UTF8]{ctex}  % 支持中文
%\title{我的份}
%\begin{document}
%\maketitle
\title{
	\includegraphics[scale = 1]{HDU.png}\\
    \vspace{1in}
    \textmd{ \Huge\textbf{成长社会学}}\\
    \textmd{\textbf{期末心得报告}}\\
   	\vspace{3.5in}
   	\large{上课时间：周一 10、11}\\
	\textmd{姓名：罗汉东}\\
	\textmd{学号：19041822}\\
	\textmd{指导老师：钱志远 }\\
}

\begin{document}
\maketitle
\newpage
我回顾了我学习《成长社会学》这门课的心路历程，
在本学期的前期，老师就自己的人生经历，分享了自己是如何进行阶层跃迁的，自己是如何从一个农村小伙，一步步网上爬最终成为博士以及大学讲师的故事，我的心潮澎湃，我觉得自己应该向老师学习。其中老师所讲的自己和自己女友的爱情故事则更加引人入胜，这让我对爱情产生的思考，在我看来爱情应该是两个相爱的人互相扶持、互相守护、共同经历人生所经历的一切，而不是单方面的追求。彼此相互了解和理解，共同经历人生中的种种困难和挫折，共同面对人生中的各种风雨和坎坷，共同面对人间的冷暖。爱情可以是亲密相伴的陪伴，也可以是分分合合的依靠。

爱情是需要经营的，经营的好，可以让彼此感情得以升华，感情得以升华的同时，也能够让彼此的价值观得以升华，能够让彼此的生活更加充实，有了一个良好的生活环境，有了一个健全的人格，有了一个完善的人格体系，对人生的态度就会更加理性和客观。爱情是需要经营的，如果经营不好，即使结婚后，两个人仍然会是一对形同路人的夫妻，没有一点爱情的升华。但是爱情却可以升华，如果两个人能够相处的更加和睦、相互理解，能够共同面对人生中的风雨和坎坷。爱情是需要经营的，如果经营不好，即使结婚，婚姻的美好也将会被生活埋没。

在之后观看的的关于婚姻的综艺，老师支出了对于婚姻的稳定和对爱情的追求之间的矛盾关系，综艺中的夫妻们也很大程度上也在苦苦维持爱情的状态。不过，我们在观看中感触到了很多，婚姻并不是两个人的事情。如果我们不想婚姻不稳，就要把对彼此的爱和婚姻的责任看得重一点。婚姻不仅是两个人相互扶持，还涉及到两个家庭成员，一起生活，生儿育女，一起解决生活中的琐事，一起为了家庭和个人的未来奋斗。我们要努力，我们要爱自己，我们要爱对方，我们都要学会如何爱人。婚姻也是爱自己，爱家人，爱生活。如果没有家人的爱和包容和体谅，就没有一起走过的路，爱是婚姻的前提。

婚姻的成功不仅仅需要婚姻的保鲜，更需要经营。我的家庭，我们的婚姻，都需要经营，家人的陪伴也是必不可少的。在节目中，老师说婚姻不仅仅是两个人的事情，还涉及到家庭，两个家庭之间要有一种默契和包容，也需要相互尊重。婚姻中，如果不尊重自己，不尊重家庭，没有包容，没有体谅，不尊重对方，不爱自己，不爱对方，那么，婚姻一定会出现问题，出现问题了，我们要学习爱自己，也要懂得爱人。我认为，婚姻不是一个人的事情，是两个家庭的事情，我们都是为了个人的生存努力，一起共同生活，我们共同为了家庭付出努力，为了家庭而努力。如果只是为了婚姻做一些小事情，这种婚姻只能说是一个美好的结果，并不能让婚姻变得更坚固。

老师也对社会的内卷进行了讨论，所谓内卷化，指一种社会（文化模式）某一发展阶段达到某种确定的形式之后，这种形式便停滞不前，难以转化为另一种高级模式的现象，从而把自我锁死在低水平状态上，周而复始地循环。而大家都对其有自己的思考，我们的生活节奏，我们的工作与生活节奏，都在不断地被加快，而这些节奏的不断加快却带来了我们无法想象的结果"社会的内卷化。在这里，我们对于社会的内卷化也有了更深层的了解，那就是我们的生活方式被快节奏的生活方式所改变。比如，我们的生存方式被快速的物质化，从而使我们的生活质量急剧下降，而不再是一种生命质量的提升，也不再是一种健康的生活方式。

之后我们分成小组进行讨论，我们各自都向大家分享了自己在人生路上的成长瞬间，我分享了自己在学校里参加学科竞赛时做出的思考和感悟，大家一起分享了自己的成长故事，这些故事和经历都是我们成长的见证，都是成长的宝贵财富，让我感觉这一次学习的收获远远要比以前的收获大许多。最后我们将我们的各种观点进行总结，上台演讲，也得到了老师的肯定和鼓励。同学们都说，上了这堂课，我们感触颇深，老师讲得不仅是学科知识，还有对我们的影响、帮助和启发。



\end{document}